% Regression: the placements the automatic station cannot read, pinned at
% the historical south default they had before the station existed.
%
% The station reads the frozen wire records in frame coordinates: an
% endpoint resolves to a cell, and the direction its ink leaves by is the
% sign of the difference between two cell numbers.  A placement the frame
% does not chart has no cell to difference, and a route that turns where its
% own waypoints say does not leave along the chord between its ends.  For
% those the scan reserves nothing and the station falls back to south, which
% is exactly where these labels stood before the station was added: no
% figure regresses, and none of them improves either.
%
% Each picture below is a placement the frame reading cannot answer.  They
% are pinned so the reading that answers them -- bearings between resolved
% points, turned into the atom's own frame, issue 6195 -- has a target to
% flip, and so a later change cannot quietly move them without saying so.
\documentclass{standalone}
\usepackage{tenkz}
\makeatletter
\ExplSyntaxOn
\file_input:n { tenkz-kernel.code.tex }
\int_new:N \l__tenkz_test_side_count_int
\tl_new:N \l__tenkz_test_side_label_tl
\cs_new_eq:NN
  \__tenkz_test_side_auto:nN \__tenkz_kernel_r_label_auto_side:nN
\cs_set_protected:Npn \__tenkz_kernel_r_label_auto_side:nN #1#2
  {
    \__tenkz_test_side_auto:nN {#1} #2
    \int_gincr:N \l__tenkz_test_side_count_int
    \__tenkz_model_get:nnN {#1} {label} \l__tenkz_test_side_label_tl
    \str_case:VnF \l__tenkz_test_side_label_tl
      {
        {P}
          {
            \str_if_eq:VnF #2 {s}
              { \errmessage{off-cell~named~end~answered:~update~the~pin} }
          }
        {Q}
          {
            \str_if_eq:VnF #2 {s}
              { \errmessage{geometric~waypoint~answered:~update~the~pin} }
          }
        {R}
          {
            \str_if_eq:VnF #2 {s}
              { \errmessage{interior~waypoint~answered:~update~the~pin} }
          }
        {T}
          {
            \str_if_eq:VnF #2 {s}
              { \errmessage{routed~carrier~bead~answered:~update~the~pin} }
          }
      }
      { \errmessage{unexpected~label~in~gap~fixture} }
  }
\ExplSyntaxOff
\makeatother
\begin{document}
\tenkzkernel
% A named endpoint whose atom stands a step off the frame owns no cell, so
% no difference of cell numbers says which way the wire leaves it.
\begin{tenkz}[lattice={2x2}, bonds=none]
  \tn[at=(1,1), name=A]{}
  \tn[at=1 s of A, name=P, skin=dot]{P}
  \tn[at=(2,2), name=C]{}
  \tnwire{P}{C}
\end{tenkz}
% A waypoint spelt as a relative address names a point, not a cell, so the
% departure bearing is a property of the resolved point.
\begin{tenkz}[lattice={2x3}, bonds=none]
  \tn[at=(1,1), name=Q, skin=dot]{Q}
  \tn[at=(1,3), name=E]{}
  \tnwire[kind=string, via={1 s of Q}]{Q}{E}
\end{tenkz}
% An atom named as an interior waypoint is entered and left by the route,
% and neither of the wire's two ends names it.
\begin{tenkz}[lattice={3x1}, bonds=none]
  \tn[at=(1,1), name=F]{}
  \tn[at=(2,1), name=R, skin=dot]{R}
  \tn[at=(3,1), name=H]{}
  \tnwire[kind=string, via={(2,1)}]{F}{H}
\end{tenkz}
% A bead on a string that bends through its own waypoints: which faces the
% carrier crosses at the bead is a property of the route at that arc length.
\begin{tenkz}[lattice={3x3}, bonds=none]
  \tn[at=(1,1), name=J]{}
  \tn[at=(1,3), name=K]{}
  \tnwire[kind=string, name=RS, via={(3,1),(3,3)}]{J}{K}
  \tn[at=on RS 0.5, skin=dot]{T}
\end{tenkz}
\ExplSyntaxOn
\int_compare:nNnF { \l__tenkz_test_side_count_int } = {4}
  { \errmessage{gap~fixture~did~not~station~four~labels} }
\ExplSyntaxOff
\end{document}
